\documentclass[11pt]{article}

\usepackage[margin=1in]{geometry}
\usepackage{amsmath,amssymb,amsthm,mathtools}
\usepackage[hidelinks]{hyperref}

\newtheorem{theorem}{Theorem}[section]
\newtheorem{lemma}[theorem]{Lemma}
\newtheorem{corollary}[theorem]{Corollary}
\theoremstyle{remark}
\newtheorem{remark}[theorem]{Remark}

\newcommand{\F}{\mathbb F}
\newcommand{\Tr}{\operatorname{Tr}}
\newcommand{\Fr}{\operatorname{Fr}}
\newcommand{\AS}{\wp}
\newcommand{\Lam}{\Lambda}

\title{A Tower-Recursive Source for Every Gold Diagonal}
\author{a9lim}
\date{August 2026}

\begin{document}
\maketitle

\begin{abstract}
Let $m$ be a power of two, let $K_m=\F_{2^m}$ be the canonical finite
nim-field, and write $e_i=2^i$ for its Conway bit basis.  For the Gold form
\[
 Q_a(x)=\Tr_{K_m/\F_2}\!\left(x^{1+2^a}\right),
\]
the diagonal constants $Q_a(e_i)$ are independent of the polar form. We give
a recursion for the trace-dual diagonal of every scaled Gold component,
prove that the unscaled dual always descends to the half-field, and conclude
that it is an Artin--Schreier value $w^2+w$.  Thus every Gold exponent has a
single tower-constructed source $w$ satisfying
\[
 \Tr\!\left((w^2+w)e_i\right)=Q_a(e_i)\qquad(0\le i<m).
\]
The construction works uniformly for odd and even $a$. It supplies the
algebraic source used by the companion weighted-source Witt--FIFO theorem,
while leaving the strategic normal-play argument logically separate.
\end{abstract}

\section{Introduction and main theorem}

For a scale $c\in K_m$, define the scaled Gold component
\[
 Q^{(m)}_{a,c}(x)=\Tr_m\!\left(c\,x\Fr^a(x)\right),
 \qquad \Fr(x)=x^2,
\]
where $\Tr_m=\Tr_{K_m/\F_2}$.  Its canonical-basis diagonal extends uniquely
to an $\F_2$-linear functional
\[
 \ell^{(m)}_{a,c}\!\left(\sum_i x_i e_i\right)
   =\sum_i x_i Q^{(m)}_{a,c}(e_i).
\]
The absolute trace pairing on a finite separable field is nondegenerate, so
there is a unique $\Lam^{(m)}_{a,c}\in K_m$ such that
\begin{equation}\label{eq:dual}
 \Tr_m\!\left(\Lam^{(m)}_{a,c}e_i\right)
 =\Tr_m\!\left(c\,e_i\Fr^a(e_i)\right)
 \quad\text{for every }i<m.
\end{equation}
Write $\Lam^{(m)}_a=\Lam^{(m)}_{a,1}$.

\begin{theorem}[all-exponent Gold-diagonal source]\label{thm:main}
For every power of two $m\ge2$ and every $a\ge0$:
\begin{enumerate}
\item $\Lam^{(m)}_{a,c}$ is given by the quadratic-tower recursion in
Theorem~\ref{thm:recursion}, using only nim addition and multiplication,
Frobenius, relative trace, and the canonical tower generator;
\item the unscaled dual descends:
\[
 \Lam^{(m)}_a\in K_{m/2}\subset K_m;
\]
\item there is $w^{(m)}_a\in K_m$ with
\[
 \Lam^{(m)}_a=\AS(w^{(m)}_a)
   =(w^{(m)}_a)^2+w^{(m)}_a;
\]
and therefore
\[
 Q_a(e_i)=\Tr_m\!\left(\AS(w^{(m)}_a)e_i\right)
 \quad(0\le i<m).
\]
\end{enumerate}
The two possible sources are $w^{(m)}_a$ and $w^{(m)}_a+1$.
\end{theorem}

The theorem is about sourcing the diagonal constants. It does not itself
infer a $P$-set theorem or a linking strategy. Those are supplied by the
separate weighted-source Witt--FIFO theorem; this proof/semantics boundary is
essential even though the companion normal-play construction supplies the
remaining game-semantic layer.

\section{The quadratic nim-tower block}

Put $m=2M$.  The canonical nim tower has
\[
 K_m=K_M(u)=K_M\oplus uK_M,
 \qquad u=e_M,
\]
and the nontrivial $K_M$-automorphism $\sigma=\Fr^M$ satisfies
\[
 \sigma(u)=u+1.
\]
The canonical basis splits as
\[
 e_0,\ldots,e_{M-1},\quad ue_0,\ldots,ue_{M-1};
\]
indeed $ue_j=e_{M+j}$.  Let
\[
 T=T_{m/M}=1+\sigma:K_m\longrightarrow K_M
\]
be the relative trace.

\begin{lemma}[trace-pairing blocks]\label{lem:blocks}
For $A,B,e\in K_M$,
\begin{align}
 T\bigl((A+uB)e\bigr)&=Be,\label{eq:lowblock}\\
 T\bigl((A+uB)ue\bigr)&=(A+B)e.\label{eq:topblock}
\end{align}
\end{lemma}

\begin{proof}
Because $\sigma$ fixes $A,B,e$ and sends $u$ to $u+1$,
\[
 (A+uB)e+\sigma((A+uB)e)
 =(A+uB)e+(A+(u+1)B)e=Be.
\]
For the upper block, apply the same calculation to $(A+uB)u$:
\begin{align*}
 &(A+uB)u+(A+(u+1)B)(u+1)\\
 &\qquad=Au+u^2B+Au+A+(u+1)^2B=A+B,
\end{align*}
since characteristic two cancels the doubled $Au$, $u^2B$, and $uB$
terms.  Multiplication by $e$ gives~\eqref{eq:topblock}.
\end{proof}

Now write $c=c_0+uc_1$ and set
\begin{equation}\label{eq:scales}
 d_0=T(c),\qquad
 d_1=T\!\left(c\,u\Fr^a(u)\right).
\end{equation}
For a lower basis vector $e\in K_M$,
\begin{align*}
 Q^{(m)}_{a,c}(e)
 &=\Tr_M\!\left(T(c)e\Fr^a(e)\right)
 =Q^{(M)}_{a,d_0}(e).
\end{align*}
For its upper mate $ue$,
\begin{align*}
 Q^{(m)}_{a,c}(ue)
 &=\Tr_M\!\left(T(cu\Fr^a(u))e\Fr^a(e)\right)
 =Q^{(M)}_{a,d_1}(e).
\end{align*}
Thus $d_0$ and $d_1$ are exactly the smaller-field scales of the two
diagonal blocks.

\begin{theorem}[scaled dual recursion]\label{thm:recursion}
Let
\[
 L_0=\Lam^{(M)}_{a,d_0},\qquad
 L_1=\Lam^{(M)}_{a,d_1},
\]
with $d_0,d_1$ as in~\eqref{eq:scales}.  Then
\begin{equation}\label{eq:recursion}
 \boxed{\quad
 \Lam^{(2M)}_{a,c}=(L_0+L_1)+uL_0.
 \quad}
\end{equation}
The base case over $K_1=\F_2$ is $\Lam^{(1)}_{a,c}=c$.
\end{theorem}

\begin{proof}
Write the candidate in~\eqref{eq:recursion} as $A+uB$, where
$B=L_0$ and $A=L_0+L_1$.  By Lemma~\ref{lem:blocks}, its relative
trace pairing with a lower basis vector $e$ is $L_0e$, while its pairing
with $ue$ is
\[
 (A+B)e=(L_0+L_1+L_0)e=L_1e.
\]
After applying $\Tr_M$, these are exactly the two defining dualities for
the lower and upper diagonal blocks.  Nondegeneracy of the absolute trace
pairing makes the candidate unique, proving~\eqref{eq:recursion}.
\end{proof}

The scale $d_1$ may change at every tower level; the recursion transports that
change rather than requiring one fixed dual element throughout the tower.

\section{Descent and the Artin--Schreier source}

\begin{lemma}[unscaled half-field descent]\label{lem:descent}
For $m=2M$,
\[
 \Lam^{(m)}_a\in K_M.
\]
\end{lemma}

\begin{proof}
In the unscaled case $c=1$, the lower scale is
\[
 d_0=T(1)=1+1=0.
\]
Linearity (or the recursion from the base case) gives
$L_0=\Lam^{(M)}_{a,0}=0$.  Theorem~\ref{thm:recursion} therefore reduces to
\[
 \Lam^{(2M)}_a=L_1\in K_M.
\]
\end{proof}

\begin{lemma}[Artin--Schreier exact sequence]\label{lem:as}
For a finite field $K=\F_{2^m}$,
\[
 \operatorname{im}(x\mapsto x^2+x)=\ker(\Tr_{K/\F_2}).
\]
Each element of this common hyperplane has exactly two preimages, differing
by $1$.
\end{lemma}

\begin{proof}
The Artin--Schreier map $\AS(x)=x^2+x$ is $\F_2$-linear and has kernel
$\F_2=\{0,1\}$, so its image has $|K|/2$ elements.  Frobenius invariance of
trace gives
\[
 \Tr(\AS(x))=\Tr(x^2)+\Tr(x)=\Tr(x)^2+\Tr(x)=0,
\]
so the image lies in the trace-zero hyperplane.  The finite-field trace is a
nonzero $\F_2$-linear functional, hence its kernel also has $|K|/2$ elements.
The inclusion is equality.  Every nonempty fibre is a coset of the two-element
kernel.
\end{proof}

\begin{proof}[Proof of Theorem~\ref{thm:main}]
Theorem~\ref{thm:recursion} proves the constructive first assertion and
Lemma~\ref{lem:descent} proves the second.  If
$\lambda=\Lam^{(2M)}_a\in K_M$, trace transitivity gives
\[
 \Tr_{2M}(\lambda)
 =\Tr_M\!\left(\Tr_{K_{2M}/K_M}(\lambda)\right)
 =\Tr_M(2\lambda)=0.
\]
Lemma~\ref{lem:as} supplies $w$ with $w^2+w=\lambda$.  Substituting this
identity into the defining trace duality~\eqref{eq:dual} yields the stated
formula for every basis coin.  The fibre statement in Lemma~\ref{lem:as}
gives the two sources $w$ and $w+1$.
\end{proof}

\section{Implementation and formal boundary}

The Rust core implements the proof literally:
\begin{itemize}
\item \texttt{gold\_component\_diagonal\_dual(m,a,c)} evaluates
Theorem~\ref{thm:recursion};
\item \texttt{gold\_diagonal\_dual(m,a)} specializes to $c=1$; and
\item \texttt{gold\_diagonal\_artin\_schreier\_source(m,a)} uses the existing
exact finite-nim-field solver to choose one of the two sources.
\end{itemize}
The implementation is checked directly against all defining trace-pairing
equations for every scale over $K_m$ at $m\le8$.  For the unscaled theorem it
checks half-field descent, the Artin--Schreier equation, and every diagonal
coordinate at representative exponents through the full $m=128$ backend.  The
previous even-$a$ values
\[
 \Lam^{(4)}_2=0,
 \quad\Lam^{(8)}_2=6,
 \quad\Lam^{(16)}_2=102,
 \quad\Lam^{(32)}_2=24582
\]
and $\Lam^{(16)}_4=31$, $\Lam^{(32)}_4=8030$ are pinned as regression values;
their drift is now explained and handled by the scale recursion.

The Lean module \texttt{formal/Ogdoad/GoldDiagonal.lean} kernel-checks the
quadratic trace-block formulas, the dual-block reconstruction, the vanishing
of the absolute trace of a quadratic-extension base element, additivity of the
Artin--Schreier map, the one-step tower lift of a source, and the full
finite-field equivalence
$\Tr(\lambda)=0\Longleftrightarrow\lambda=w^2+w$. The last result is proved
from rank--nullity, trace surjectivity, and the two-element kernel, rather than
introduced as an axiom. The module does not encode the finite nim tower or
replace the Rust implementation with a Lean finite-field model.

\section{Integration with Gold normal play}

For every exponent, the diagonal is supplied by one deterministically selected tower
constant through nim multiplication, Frobenius, trace, and finally one square
plus XOR at the source identity $\Lam=\AS(w)$. No coordinate-by-coordinate
evaluation of $Q_a$ occurs in the recursion. Unlike the $a=1$ Fermat-coin
formula, the final $w$ is selected by exact finite-field linear algebra; the
theorem does not claim a comparably simple closed expression or that this
selection itself satisfies the strategic naturality axioms.

The strategic layer is supplied separately by weighted-source Witt--FIFO. Its
source-pair weights consume these diagonal bits locally, the matching theorem
forces the exact charge in every dimension, and the phase-aware claim lantern
gives an ordinary normal-play P-set $\{Q_a=0\}$. The present theorem supplies
the all-$a$ algebraic input to that construction; it does not collapse the
distinction between sourcing and play semantics.

\end{document}
